\subsection{\bt{Questions and setup}{问题与设置}}
\bi{The evaluation follows the layer's responsibilities. Reliance: (Q1) do maintained definitions stay correct under data-only changes, independently of what agents learned? (Q2) Does use-time enforcement hold under concurrent writes, and at what cost? Maintenance: (Q3) when does repair restore a broken definition, and when must it invalidate it? (Q4) What does maintenance cost? Value: (Q5) what does maintained shared memory do for LLM agents, against exploring alone and retrieving examples? Publication: (Q6) what does admission refuse, and how much faster does the memory make consumers? Improvement: (Q7) do consumers adopt optimization revisions, and does the layer fall back when an added premise breaks? Lifecycle: (Q8) does one library keep serving arriving consumers as updates accumulate? Q1--Q4 are answered at the system level without LLMs; Q5--Q8 use complete agents, and Q7 adds a system-level test. The machine has 48 Hygon cores and 98\,GiB RAM; PostgreSQL 18.6 runs in disposable databases with WAL and fsync disabled. The synthetic retail fixture uses TPC-DS-style column names; one million store-sales rows come with about 396,000 returns and 500,000 catalog sales. The library replay also runs on three schemas not written for this paper: TPC-H SF1 (date columns, snowflake joins), SSB SF1 (yyyymmdd date keys, other calendar columns), and the BIRD financial database~\cite{li2023bird} of a real bank (repeated column names, one-to-one joins); their libraries come from each benchmark's own queries, and one generator instantiates the eleven changes from a short description of each schema. Table~\ref{tab:ladder} lists the compared methods; revalidation methods share conditions, repair, regression test, and request coalescing. The regression test compares with the agent's own learning SQL as of learning time; only the full recheck uses the benchmark's gold SQL end to end, where \system\ with that reference answers the same as \system.}
{评估按这一层承担的职责展开。依赖：（Q1）维护后的定义能否在纯数据变化下保持正确，且不受智能体学到什么的影响？（Q2）存在并发写入时，使用时的强制保证是否成立，代价如何？维护：（Q3）修复何时能恢复被破坏的定义，何时必须使定义失效？（Q4）维护代价如何？价值：（Q5）与独自探索和检索示例相比，维护后的共享记忆对大模型智能体有何作用？发布：（Q6）准入拒绝了什么，记忆让使用者快了多少？改进：（Q7）使用者是否采用优化修订，新增前提被破坏时这一层能否回退？生命周期：（Q8）随着更新累积，一个定义库能否持续服务陆续到来的使用者？Q1--Q4 在系统层回答，不调用大模型；Q5--Q8 使用完整的智能体，Q7 另有一个系统层测试。实验机有 48 核 Hygon 和 98\,GiB 内存；PostgreSQL 18.6 在可丢弃的数据库中运行，关闭 WAL 与 fsync。合成零售数据采用 TPC-DS 风格列名；100 万行门店销售对应约 39.6 万行退货和 50 万行目录销售。定义库回放还在三种不是为本文编写的模式上运行：TPC-H SF1（日期列、雪花型连接）、SSB SF1（yyyymmdd 日期键、另有名称的日历列）和 BIRD 的 financial 库~\cite{li2023bird}，一家银行的真实数据（列名重复、一对一连接）；各定义库由该基准自己的查询导出，同一个生成器按每种模式的简短描述生成全部 11 种变化。表~\ref{tab:ladder} 列出被比较的方法；各重验证方法共用条件、修复、回归测试与请求合并。回归测试以智能体自己的学习 SQL 按学习时刻比较；只有整定义重查在端到端实验中使用基准的标准答案 SQL，而以该参照的 \system\ 在那里与 \system\ 答对的题数相同。}

\begin{table}[t]
\caption{\bt{Maintenance strategies for shared definitions and the studies that compare them: R, library replay; C, controlled cost benchmark; E, end-to-end scenarios (Section~\ref{sec:evaluation}). ``Write during use'': a write that commits between validation and the query that relies on it.}{共享定义的维护方式及比较它们的实验：R，定义库回放；C，受控代价实验；E，端到端情形（第~\ref{sec:evaluation} 节）。“使用期间写入”：在验证与依赖它的查询之间提交的写入。}}
\label{tab:ladder}
\footnotesize
\setlength{\tabcolsep}{3pt}
\begin{tabularx}{\linewidth}{@{}llLLLL@{}}
\toprule
\bt{Strategy}{方式} & \bt{Used in}{实验} & \bt{Benign append}{正常追加} & \bt{Data-only break}{纯数据破坏} & \bt{Write during use}{使用期间写入} & \bt{Work per update}{每次更新的工作}\\
\midrule
\swatch{mSchema}\ \ifbilingual\bhl{Schema-change\\invalidation}{按模式变更失效}\else\bhl{Schema-change\\invalidation}{}\fi & R, E & \bt{keep}{保留} & \bt{stale use}{过期使用} & \bt{unchecked}{不检查} & \bt{none}{无}\\
\swatch{mTable}\ \bhl{Table tests}{表级测试} & R & \bt{keep}{保留} & \bt{invalidate}{失效} & \bt{check, then execute}{先检查后执行} & \bt{declared tests per table}{按表声明的测试}\\
\swatch{mRevoke}\ \ifbilingual\bhl{Invalidate-\\on-write}{写入即失效}\else\bhl{Invalidate-\\on-write}{}\fi & R, E & \bt{relearn}{重学} & \bt{relearn}{重学} & \bt{check, then execute}{先检查后执行} & \bt{LLM relearning}{大模型重学}\\
\swatch{mDef}\ \bhl{Definition-level}{定义级} & R, C, E & \bt{keep}{保留} & \bt{repair}{修复} & \bt{check, then execute}{先检查后执行} & \bt{conditions of affected definitions}{受影响定义的条件}\\
\swatch{mCache}\ \ifbilingual\bhl{Check-result\\cache}{检查结果缓存}\else\bhl{Check-result\\cache}{}\fi & R, C & \bt{keep}{保留} & \bt{repair}{修复} & \bt{check, then execute}{先检查后执行} & \bt{distinct changed checks}{变化的不同检查}\\
\swatch{mCond}\ \system & R, C, E & \bt{keep}{保留} & \bt{bounded repair}{有界修复} & \bt{same snapshot}{同一快照} & \bt{distinct changed conditions}{变化的不同条件}\\
\bottomrule
\end{tabularx}
\end{table}

